\capitulo{7}{Conclusiones y Líneas de trabajo futuras}

El proyecto \nombreProyecto ha permitido diseñar e implementar una solución
completa para el análisis temático cualitativo de memorias de \ACRshort{tfg}
del grado en Ingeniería Informática de la Universidad de Burgos publicadas en
GitHub, cubriendo todo el ciclo de vida del dato: ingesta, procesamiento,
modelado y visualización.

A nivel de objetivos, se considera que el proyecto se ha cumplido de forma
satisfactoria. Se ha construido un \gls{pipeline} modular con tres bloques bien
diferenciados (\textit{backend-ingestion}, \textit{backend-processing} y
\textit{frontend-web}), con una arquitectura hexagonal en backend que facilita
la mantenibilidad, la escalabilidad y la sustitución de componentes. Además, se
ha definido un contrato de datos estable en formato Parquet, lo que mejora la
trazabilidad experimental y desacopla claramente la generación de resultados de
su consumo en la interfaz.

Desde el punto de vista experimental, se han entrenado y comparado cuatro
enfoques de modelado de temas (\ACRshort{lda}, BERTopic, Top2Vec y FASTopic)
sobre cuatro fuentes textuales (\textit{abstracts}, \textit{issues},
\textit{readmes} y cuerpo de las memorias). En el estado actual del proyecto,
BERTopic ofrece la mejor coherencia en \textit{issues}, \textit{readmes} y
cuerpo de las memorias, mientras que Top2Vec destaca en \textit{abstracts}. Por
su parte, LDA presenta un comportamiento competitivo en varios escenarios con
un coste computacional significativamente menor. Este resultado evidencia un
compromiso técnico relevante entre calidad semántica e inversión de recursos.

Estos resultados confirman que no existe un único modelo óptimo universal, sino
que el rendimiento depende del tipo de texto analizado y de su granularidad
semántica. También se observa un compromiso claro entre calidad temática y
coste computacional: los modelos basados en \textit{embeddings} suelen mejorar
la coherencia en varios escenarios, a costa de mayores tiempos de ejecución.

Como valoración crítica, identificamos tres limitaciones principales. En primer
lugar, la evaluación se apoya mayoritariamente en métricas automáticas de
coherencia, que no sustituyen una validación cualitativa experta. En segundo
lugar, la heterogeneidad de los repositorios de origen introduce ruido
documental y estructural que afecta a la estabilidad del modelado. En tercer
lugar, algunos modelos basados en \textit{embeddings} presentan costes de
ejecución elevados, lo que condiciona la escalabilidad en corpus más grandes.

Además de los resultados cuantitativos, uno de los aportes principales del
trabajo es haber construido una base software reutilizable: el sistema permite
reejecutar fases de forma independiente, comparar configuraciones de manera
homogénea y exponer resultados en una interfaz interactiva orientada a análisis
exploratorio.

\section{Líneas de trabajo futuras}

Como evolución natural del proyecto, se plantean las siguientes líneas:

\begin{itemize}
      \item Mejora de la interpretación de los temas generados. Actualmente, los temas
            obtenidos por los modelos se representan mediante distribuciones de
            probabilidad sobre conjuntos de términos, lo que requiere una interpretación
            posterior por parte del usuario para identificar de manera precisa la temática
            asociada. Habitualmente, esta tarea se realiza mediante validación humana. Como
            línea de trabajo futura, se plantea incorporar modelos grandes de lenguaje
            (\acrfull{llm}) para automatizar el proceso de etiquetado e interpretación
            semántica de los temas, ya que las pruebas preliminares realizadas han mostrado
            resultados satisfactorios en tareas de clasificación e identificación temática.
      \item Ampliar la evaluación con más métricas automáticas (por ejemplo, diversidad
            temática y estabilidad entre ejecuciones) y reforzar la validación humana con
            protocolos estandarizados de evaluación.
      \item Incorporar análisis temporal para estudiar la evolución de temáticas por curso
            académico y detectar tendencias emergentes.
      \item Mejorar la normalización lingüística en español (tratamiento de tecnicismos,
            acrónimos y expresiones de dominio) para aumentar la calidad semántica de los
            temas.
      \item Explorar variantes de modelado adicionales (modelos jerárquicos, dinámicos o
            nuevas configuraciones de \textit{sentence-transformers}).
      \item Optimizar rendimiento mediante caché de \textit{embeddings}, ejecución paralela
            y otras estrategias de reducción de coste computacional.
      \item Mejorar la interfaz de usuario: filtros académicos más ricos, exportación de
            informes para toma de decisiones y mejoras en accesibilidad para usuarios con
            necesidades especiales.
\end{itemize}

En definitiva, este \ACRshort{tfg} no solo entrega un prototipo funcional, sino
una base técnica sólida y extensible sobre la que continuar investigación
aplicada en analítica temática de documentación académica.
